\chapter{Anforderungsanalyse}
\label{chap:anforderungsanalyse}

Dieses Kapitel leitet aus der Zielsetzung der Arbeit
(Abschnitt~\ref{sec:zielsetzung}) die Anforderungen an das zu entwickelnde System ab.
Ausgangspunkt ist das Einsatzszenario der temporären Bühnen- und
Veranstaltungsbeleuchtung, aus dem sich zunächst die Betriebsmodi ergeben. Auf dieser
Grundlage werden die funktionalen Anforderungen komponentenbezogen für Bridge,
Leuchtröhren und Kommunikation formuliert und um nicht-funktionale
Qualitätsanforderungen ergänzt. Den Abschluss bildet die Abgrenzung jener Aspekte, die
nicht Gegenstand der Arbeit sind. Die hier festgelegten Anforderungen bilden den
Maßstab, an dem das System in Kapitel~\ref{chap:evaluierung} gemessen wird.

\section{Systemkontext und Einsatzszenarien}
\label{sec:systemkontext}

Das System richtet sich an die Beleuchtung von Bühnen und Veranstaltungen, insbesondere
an temporäre Aufbauten, die wiederholt auf- und abgebaut werden. Kennzeichnend für
dieses Umfeld sind wechselnde Örtlichkeiten, knappe Aufbauzeiten und eine
Signalverkabelung, die bei jedem Aufbau neu verlegt werden muss
(Abschnitt~\ref{sec:problemstellung}).

Am Betrieb beteiligt sind vier Akteure. Die Lichttechnikerin oder der Lichttechniker
bedient eine bestehende Lichtsteuerung eine DMX-Konsole oder eine
Steuerungssoftware, die den Lichtzustand als DMX-Universum ausgibt. Die Bridge nimmt
dieses Signal entgegen und verteilt die für die Röhren bestimmten Kanäle. Die
Leuchtröhren setzen die empfangenen Werte in Licht um. Konfiguriert werden Bridge und
Röhren unmittelbar am Gerät, ohne zusätzliche Software.

Aus diesem Kontext ergeben sich drei typische Einsatzszenarien. Im ersten steht eine
vollständige Lichtsteuerung zur Verfügung. die Röhren sollen sich wie gewöhnliche
DMX-Geräte in die bestehende Anlage einfügen, jedoch ohne Signalkabel auskommen. Im
zweiten fehlt eine Lichtsteuerung, etwa bei kleineren Veranstaltungen. Das System soll
dann selbstständig ein zusammenhängendes Lichtbild erzeugen. Im dritten wird eine
einzelne Röhre ohne weitere Infrastruktur als reine Lichtquelle eingesetzt. Diese drei
Szenarien führen unmittelbar auf die im folgenden Abschnitt festgelegten
Betriebsmodi.

\section{Betriebsmodi}
\label{sec:betriebsmodi}

Aus den beschriebenen Einsatzszenarien ergeben sich drei Betriebsmodi: Der
\textit{Standalone-Modus} bezeichnet den Betrieb einer Leuchtröhre ohne Verbindung zur
Bridge, der \textit{Master-Modus} den eigenständigen Betrieb der Bridge ohne externe
Steuersignale und der \textit{DMX-Modus} den Funkempfang der Leuchtröhre, während die
Bridge das DMX-Signal einer externen Lichtsteuerung entgegennimmt.

\subsection{Standalone-Modus (Leuchtröhre)}
\label{subsec:standalone-modus}

Im Standalone-Modus kann die LED-Leuchtröhre vollständig
autark betrieben werden, ohne dass eine Verbindung zur Bridge erforderlich ist. 
Dieser Betriebsmodus bietet eine hohe Flexibilität für Anwendungsfälle, 
in denen lediglich eine einfache Lichtquelle ohne komplexe Steuerung benötigt wird.

Über das Menü am Gerät wählt der Benutzer die gewünschte Farbe sowie einen der
vorgespeicherten Effekte aus beide werden unmittelbar auf den LED-Strip angewendet.

Im Standalone-Modus entfällt jegliche Kommunikation mit externen Geräten. Die
Leuchtröhre arbeitet vollständig eigenständig und benötigt lediglich eine
Stromversorgung.

\subsection{Master-Modus (Bridge)}
\label{subsec:master-modus}

Im Master-Modus empfängt die Bridge keine Steuersignale von außen, sondern erzeugt den
Lichtzustand selbst. Sie durchläuft dabei selbstständig eine Folge vorgegebener Farben
und Effekte und überträgt das Ergebnis über die Funkschnittstelle an alle Leuchtröhren.
Das Tempo dieses Ablaufs lässt sich am Bedienmenü der Bridge über die Schlagzahl
(Beats per Minute) einstellen, sodass sich die Effektwechsel an die Musik einer
Veranstaltung anpassen lassen.

Da sämtliche Röhren denselben Datenstrom empfangen, entsteht auch ohne externe
Lichtsteuerung ein kohärentes Gesamtbild über alle Geräte hinweg. Der Modus eignet sich
damit für einfachere Lichtinstallationen sowie für Veranstaltungen, bei denen keine
Lichtsteuerung zur Verfügung steht oder deren Einrichtung nicht gerechtfertigt wäre.

\subsection{DMX-Modus (Bridge und Leuchtröhre)}
\label{subsec:dmx-modus}

Im DMX-Modus empfängt die Bridge DMX-Steuersignale über eine
kabelgebundene DMX-Leitung und überträgt die enthaltenen Steuerinformationen
anschließend über eine Funkschnittstelle an die Leuchtröhren. Die Bridge
bezieht dabei die in den DMX-Daten enthaltenen Kanalinformationen auf die an
jeder Leuchtröhre individuell konfigurierbaren DMX-Adressen. Dadurch ist es
möglich, an jeder Leuchtröhre dieselbe DMX-Adresse einzustellen, die auch in
der verwendeten Lichtsteuerungssoftware für das jeweilige Gerät vergeben wurde.
Die Leuchtröhren empfangen die von der Bridge gesendeten DMX-Pakete und setzen die
darin enthaltenen Kanalwerte unmittelbar in entsprechende RGB-Lichtsignale um.
Jede Leuchtröhre belegt in diesem Modus 4 DMX-Kanäle. 1 Kanal zur Auswahl 
vorgespeicherter Effekte und 3 Kanäle für die RGB-Farbsteuerung. 
Eine pixelgenaue Einzelansteuerung der LEDs ist in diesem Modus nicht
vorgesehen.

Der DMX-Modus eignet sich insbesondere für Einsatzszenarien, 
in denen eine flexible Positionierung der Leuchtröhren ohne den Aufwand einer 
aufwendigen Signalverkabelung erforderlich ist. Durch den Wegfall der 
Kabelverbindungen zwischen Bridge und Leuchtröhren lässt sich der Aufbau 
erheblich beschleunigen und vereinfachen.

Über diese drei Modi hinaus ist ein kabelgebundener Betriebsmodus angestrebt, in dem
die Bridge die einzelnen Pixel der Leuchtröhren über eine Datenleitung ansteuert und
damit eine deutlich feinere Lichtgestaltung ermöglicht. Dieser Modus ist Gegenstand
einer Weiterentwicklung und nicht Teil der vorliegenden Arbeit
(Abschnitt~\ref{sec:abgrenzung}).

\section{Funktionale Anforderungen}
\label{sec:funktionale-anforderungen}

Die funktionalen Anforderungen beschreiben die konkreten Fähigkeiten und
Verhaltensweisen, die das System bereitstellen muss, um die in
Abschnitt~\ref{sec:betriebsmodi} definierten Betriebsmodi zu realisieren.
Die Anforderungen werden komponentenbezogen formuliert, um eine direkte
Grundlage für die spätere Systemarchitektur zu schaffen.

\subsection{Anforderungen an die Bridge}
\label{subsec:fa-bridge}

Die Bridge bildet das zentrale Steuerungselement des Systems und muss
dementsprechend sämtliche eingehende Signalquellen verarbeiten sowie
die erzeugten Steuerdaten an die Leuchtröhren weiterleiten können.

Im DMX-Modus muss die Bridge in der Lage sein,
DMX-Steuersignale über eine kabelgebundene DMX-Leitung gemäß dem DMX512-Standard
zu empfangen und die enthaltenen Kanalinformationen zu extrahieren.
Äquivalent dazu sollen auch Art-Net-Signale über eine Ethernet-Schnittstelle
empfangen und verarbeitet werden können.
Anschließend sind die eingehenden Signale in ein einheitliches internes
Format zu überführen, das eine modusunabhängige Weiterverarbeitung ermöglicht.
Die aufbereiteten Steuerdaten müssen anschließend gleichzeitig an alle
angeschlossenen Leuchtröhren übermittelt werden.

Zur Konfiguration des Betriebsmodus sowie weiterer systemrelevanter
Parameter muss die Bridge über ein Bedienmenü verfügen, das eine
direkte Einstellung am Gerät ohne externe Hilfsmittel ermöglicht.

\subsection{Anforderungen an die LED-Leuchtröhren}
\label{subsec:fa-roehre}

Die LED-Leuchtröhren stellen die Ausgangskomponenten des Systems dar und
müssen abhängig vom aktiven Betriebsmodus unterschiedliche Funktionalitäten
bereitstellen.

Im Standalone-Modus muss die Leuchtröhre vollständig autark operieren, ohne
dass eine Kommunikationsverbindung zu anderen Systemkomponenten erforderlich
ist. Über ein lokales Bedienmenü muss der Benutzer die gewünschte Lichtfarbe
und Animation direkt am Gerät einstellen können.

Im DMX-Modus muss die Leuchtröhre Steuerpakete über die
Funkschnittstelle empfangen und die übertragenen Kanalwerte unmittelbar in
entsprechende RGB-Lichtsignale umsetzen. Jede Leuchtröhre belegt dabei
vier DMX-Kanäle: drei Kanäle für die individuelle Steuerung der
Farbkomponenten Rot, Grün und Blau sowie einen weiteren Kanal zur Auswahl
vorgespeicherter Pixelanimationen.

Unabhängig vom Betriebsmodus muss jede Leuchtröhre über ein Bedienmenü
verfügen, über das der gewünschte Betriebsmodus sowie modusabhängige
Parameter konfiguriert werden können. Sämtliche vorgenommenen Einstellungen
müssen persistent gespeichert und beim nächsten Einschalten des Geräts
automatisch wiederhergestellt werden, damit die Geräte nach einem temporären
Stromausfall ohne erneute Konfiguration weiterarbeiten.

\subsection{Anforderungen an die Kommunikation}
\label{subsec:fa-kommunikation}

Die Kommunikationskomponente des Systems muss die zuverlässige
Signalweiterleitung von der Bridge zu den Leuchtröhren gewährleisten.

Für den drahtlosen Betriebsmodus ist eine Funkverbindung erforderlich, die
eine ausreichende Reichweite besitzt, um typische Bühnenszenarien abzudecken.
Hierbei werden 50 Meter als Mindestreichweite für den zuverlässigen Betrieb 
der Leuchtröhren als Mindestanforderung festgelegt.
Die Übertragung muss auch in hochfrequenzbelasteten Umgebungen, wie sie bei
Veranstaltungen mit parallelem Betrieb zahlreicher Funkgeräte und WLAN-Systeme
auftreten können, zuverlässig funktionieren. Die erreichbare Datenrate muss
dabei ausreichen, um die DMX-Steuerdaten für alle angeschlossenen Leuchtröhren
mit der für eine flüssige Lichtsteuerung notwendigen Aktualisierungsrate zu
übertragen.

Als Ausbaugröße wird festgelegt, dass das System bis zu 20 Leuchtröhren gleichzeitig
versorgen können muss. Da jede Röhre vier DMX-Kanäle belegt, ergibt sich daraus ein zu
übertragender Adressbereich von 80 Kanälen. Dieser Wert bildet zugleich die Grundlage
für die Auslegung der Übertragung und für die in Kapitel~\ref{chap:evaluierung}
verwendete Paketgröße.

\section{Nicht-funktionale Anforderungen}
\label{sec:nicht-funktionale-anforderungen}

Die nicht-funktionalen Anforderungen definieren Qualitätseigenschaften des
Systems, die über die reine Funktionalität hinausgehen und maßgeblich über
die praktische Einsetzbarkeit sowie die Akzeptanz des Systems entscheiden.

\subsection{Echtzeitverhalten}
\label{subsec:echtzeitverhalten}

Da das System zur Steuerung von Bühnenbeleuchtung eingesetzt wird, unterliegt
es zeitkritischen Anforderungen. Die Gesamtlatenz zwischen dem Eingang eines
Steuersignals an der Bridge und der entsprechenden Lichtausgabe an den
Leuchtröhren muss unterhalb der Wahrnehmungsschwelle des menschlichen Auges
liegen, sodass Steuereingriffe für das Publikum als unmittelbar erscheinen.

Der Critical Flicker Fusion Threshold (CFFT) ist ein psychophysikalisches Maß 
für die visuelle Zeitauflösung, also die schnellste Rate, bei der das visuelle 
System des Menschen noch einzelne Lichtsignale unterscheiden kann. Dieser Wert 
unterscheidet sich beim Menschen individuell, liegt aber typischerweise im 
Bereich von 40 bis 60 Hz~\cite{haarlem2024speed}. 

Da die einzelnen LED-Pixel nur bei neuen Instruktionen durch den Mikrocontroller
aktualisiert werden und ansonsten dauerhaft leuchten, kann dieser Grenzwert automatisch
als erfüllt angesehen werden. 

Kritischer ist jedoch der Grenzwert für die Bewegungswahrnehmung (Motion Perception Threshold), 
der bei etwa 15\,Hz liegt, optimalerweise sogar bei 24\,Hz~\cite{fujii2018}.

Einen Sonderfall stellt das typische Einsatzszenario von Bühnenbeleuchtung dar, da hier häufig 
sehr dunkle Umgebungen vorliegen. Unter solchen Bedingungen verschiebt sich die 
Flimmerverschmelzungsfrequenz aufgrund der Anatomie des menschlichen Auges nach unten.

Die Netzhaut verfügt über zwei Arten von Photorezeptoren, die sich in ihrer Funktion 
deutlich unterscheiden. Stäbchen (engl. \textit{rods}) sind hochempfindlich und ermöglichen 
das Sehen bei geringer Lichtintensität, besitzen jedoch eine vergleichsweise langsame 
Zeitauflösung, ihre Flimmerverschmelzungsfrequenz liegt bei etwa 15\,Hz. Sie dominieren 
daher das Sehen in dunklen Umgebungen. Zapfen (engl. \textit{cones}) sind hingegen für 
das Farb- und Tagsehen bei höherer Lichtintensität zuständig und erreichen mit einer 
Flimmerverschmelzungsfrequenz von bis zu 60\,Hz eine wesentlich höhere
Zeitauflösung.

Da bei Dunkelheit überwiegend die langsameren Stäbchen aktiv sind, sinkt die wahrnehmbare 
Flimmerverschmelzungsfrequenz entsprechend ab~\cite{ramunae2025}.

Ein Wert von 15\,Hz kann deshalb als genügend angesehen werden. Um
Sicherheitsreserven vorzuhalten, wird dennoch ein Grenzwert von 24\,Hz angestrebt, 
um auch in Situationen mit höherer Umgebungshelligkeit eine flüssige und reaktive 
Lichtsteuerung zu gewährleisten.

Das DMX512-Protokoll sieht eine maximale Aktualisierungsrate von 44~Hz
vor \cite{ANSI_E1.11_2024,yang2015real}. Dieser Wert liegt bereits oberhalb der angestrebten Grenzwerte
für die Bewegungswahrnehmung, sodass die Einhaltung der Echtzeit-Anforderung 
mit DMX als unkritisch angesehen werden kann.

Auch die zeitliche Konstanz der Frames ist von Bedeutung. Starke Schwankungen der
Übertragungsverzögerung (Jitter) können zu einem unregelmäßigen Lichtverhalten führen,
das selbst bei geringer bis mittlerer Bildrate als störend wahrgenommen wird.

Aus diesen Überlegungen ergeben sich drei prüfbare Anforderungen. Erstens muss die
Lichtausgabe mit mindestens 24\,Hz aktualisiert werden, ein einmal gesetzter
Lichtzustand darf im Regelbetrieb nicht länger als 41,7\,ms gehalten werden. Zweitens
darf die Gesamtlatenz zwischen dem Eingang eines Steuersignals an der Bridge und der
Lichtausgabe an der Röhre denselben Wert von 41,7\,ms nicht überschreiten. Drittens
muss der Jitter deutlich unterhalb des Abstands zweier aufeinanderfolgender Frames
bleiben, damit sich die Schwankung nicht zu einer wahrnehmbaren Unregelmäßigkeit
aufsummiert. Diese Grenzwerte bilden den Maßstab für die Messreihen in
Kapitel~\ref{chap:evaluierung}.

\subsection{Zuverlässigkeit und Robustheit}
\label{subsec:zuverlaessigkeit}

Im Veranstaltungsbetrieb beeinträchtigen Systemausfälle den Ablauf unmittelbar das
System muss daher eine hohe Betriebszuverlässigkeit gewährleisten.
Einzelne verlorene Pakete müssen ohne wahrnehmbare Wirkung auf die Lichtausgabe bleiben,
und Unterbrechungen, welche die in Abschnitt~\ref{subsec:echtzeitverhalten} genannte
Grenze von 41,7\,ms überschreiten, dürfen im Regelbetrieb nur vereinzelt auftreten.

Die drahtlose Verbindung muss auch in hochfrequenzbelasteten Umgebungen stabil
bleiben, wie sie bei Veranstaltungen mit einer Vielzahl gleichzeitig betriebener
WLAN- und Bluetooth-Geräte typischerweise auftreten. Bei einem Verbindungsverlust
muss das System ein definiertes Verhalten zeigen, wie das Einfrieren
des zuletzt empfangenen Lichtzustands, um unkontrollierte Lichtausfälle zu
vermeiden.

Darüber hinaus muss das System gegenüber Fehlbedienungen durch den Anwender
robust sein. Hierzu zählen beispielsweise die Eingabe ungültiger DMX-Adressen
oder eine fehlerhafte Moduswahl. Solche Eingaben dürfen nicht zu einem
undefinierten Systemzustand führen. Auch auf der Hardwareebene muss eine
ausreichende mechanische Stabilität gewährleistet sein, da die Geräte im
Veranstaltungsbetrieb gelegentlichen Stößen und raueren Handhabungsbedingungen
ausgesetzt sein können.

\subsection{Elektrische Belastbarkeit}
\label{subsec:elektrische-belastbarkeit}

Die elektrische Auslegung des Systems orientiert sich an der maximalen
Leistungsaufnahme der verbauten LED-Strips. Jede Leuchtröhre enthält
einen WS2815-LED-Strip mit einer Pixeldichte von 60~Pixeln pro Meter
bei einer Betriebsspannung von 12\,V. Da jeder Pixel bei vollständiger
Aussteuerung aller drei Farbkanäle eine Leistung von 0{,}2\,W aufnimmt,
ergibt sich für eine Röhrenlänge von 2\,m eine maximale Leistungsaufnahme
von
\[
  P_{\max} = 60 \times 2 \times 0{,}2\,\text{W} = 24\,\text{W}.
\]
Bei einer Versorgungsspannung von 12\,V folgt daraus ein maximaler
Betriebsstrom von
\[
  I_{\max} = \frac{P_{\max}}{U} = \frac{24\,\text{W}}{12\,\text{V}} = 2\,\text{A}.
\]

Alle stromführenden Komponenten, insbesondere Leiterbahnen, Steckverbinder
und Schutzschaltungen, müssen so dimensioniert sein, dass sie diesen Strom im
Dauerbetrieb ohne unzulässige Erwärmung oder Ausfälle führen können, die Auslegung
erfolgt dabei mit Reserve gegenüber dem berechneten Wert. Gleiches gilt für das eingesetzte Netzteil, das die geforderte
Leistung dauerhaft und ohne signifikante Spannungseinbrüche bereitstellen
muss, um einen stabilen Betrieb auch bei längeren Veranstaltungen zu
gewährleisten.

\subsection{Benutzerfreundlichkeit}
\label{subsec:benutzerfreundlichkeit}

Ein wesentliches Ziel der Systementwicklung ist die Minimierung des
Einrichtungsaufwands. Als Kriterium wird festgelegt, dass sich sämtliche für den Betrieb
erforderlichen Parameter Betriebsmodus, DMX-Adresse, Funkkanal, Farbe und Effekt
ausschließlich am Gerät einstellen lassen müssen, ohne zusätzliche Software, ohne
Netzwerkzugang und ohne schriftliche Anleitung. Die Bedienoberfläche muss dafür
selbsterklärend gestaltet sein.

Die Konfiguration von Betriebsmodus und weiteren betriebsrelevanten Parametern
erfolgt über das Bedienmenü der jeweiligen Geräte. Um die Übersichtlichkeit
zu wahren, sollen dabei ausschließlich die für den jeweils aktiven Betriebsmodus
relevanten Einstellungen angezeigt werden. Sämtliche vorgenommenen
Konfigurationen müssen persistent im Gerät gespeichert werden und stehen nach
einem Neustart ohne erneuten Einrichtungsaufwand sofort zur Verfügung.

\subsection{Wirtschaftlichkeit und Portabilität}
\label{subsec:wirtschaftlichkeit}

Das entwickelte System soll eine kosteneffiziente Alternative zu kommerziellen
drahtlosen DMX-Lösungen darstellen, die für private Personen, kleinere Produktionen und
experimentelle Installationen zugänglich ist. Bei der Komponentenauswahl
ist daher auf ein günstiges Preis-Leistungs-Verhältnis zu achten, ohne
die funktionalen und qualitativen Anforderungen zu kompromittieren. Die Materialkosten
sind damit ein Entwurfskriterium. eine betriebswirtschaftliche Bewertung des Systems
findet dagegen nicht statt (Abschnitt~\ref{sec:abgrenzung}).

Hinsichtlich der Portabilität sind eine kompakte Bauform und ein geringes
Gewicht aller Systemkomponenten anzustreben, um den Transport und den Aufbau
an wechselnden Veranstaltungsorten zu erleichtern. Darüber hinaus muss das
System ohne Anpassungen in bestehende DMX-basierte Lichtsteuerungsinfrastrukturen
integrierbar sein, um den Einsatz neben vorhandener Technik zu ermöglichen.

\section{Abgrenzung}
\label{sec:abgrenzung}

Die folgende Abgrenzung definiert, welche Funktionalitäten und Aspekte nicht Teil dieser Arbeit sind:

\begin{itemize}
    \item \textbf{DMX-Konsole}: Das entwickelte System stellt keine vollständige DMX-Konsole dar, sondern dient ausschließlich der Weiterleitung und drahtlosen Übertragung von DMX-Signalen.
    
    \item \textbf{Lichtplanung und Programmierung}: Die Planung von Lichtszenen und die Programmierung von Lichteffekten erfolgt über externe Software oder DMX-Konsolen. Das System bietet keine eigene Programmierumgebung.
    
    \item \textbf{Gerätetypen}: Das System ist auf die Ansteuerung von LED-Leuchtröhren beschränkt. Andere Beleuchtungsgeräte wie Moving Heads, Scanner oder herkömmliche PAR-Scheinwerfer werden nicht unterstützt.
    
    \item \textbf{Bidirektionale Kommunikation}: Eine bidirektionale Kommunikation ist nicht vorgesehen. Das System arbeitet ausschließlich mit unidirektionaler Signalübertragung.
    
    \item \textbf{Weitere Protokolle}: Neben DMX und Art-Net werden keine weiteren Steuerungsprotokolle wie sACN (Streaming ACN) oder KiNET unterstützt.

    \item \textbf{Kabelgebundene Pixelsteuerung}: Ein Betriebsmodus, in dem die Bridge die
    einzelnen Pixel der Leuchtröhren über eine Datenleitung ansteuert, wird als
    Weiterentwicklung angestrebt, ist jedoch nicht Gegenstand dieser Arbeit.

    \item \textbf{Wirtschaftlichkeit}: Die Materialkosten der Komponenten sind ein
    Entwurfskriterium (Abschnitt~\ref{subsec:wirtschaftlichkeit}), eine
    betriebswirtschaftliche Bewertung findet dagegen nicht statt. Das Projekt zielt nicht
    darauf ab, ein gewinnbringendes Produkt zu entwerfen, sondern wird in einem
    semi-professionellen Rahmen eigengenutzt.
\end{itemize}